\noindent\textbf{Elección del procedimiento.} Continuidad da $V_1=(A_2/A_1)V_2$. Bernoulli horizontal da $\Delta p=\rho[V_2^2-V_1^2]/2$. Al sustituir la primera relación, aparece $1-(A_2/A_1)^2=1-(D_2/D_1)^4$.

\begin{enumerate}

\item Calcular el área de entrada.

\[\fitmath{A_{1}=\frac{\pi\cdot {D_{1}}^{2}}{4}}\]

\[\fitmath{A_{1}=\frac{\pi\cdot {(0.1)}^{2}}{4}\approx 0.007853982\,\mathrm{m^{2}}}\]

\item Calcular la de garganta.

\[\fitmath{A_{2}=\frac{\pi\cdot {D_{2}}^{2}}{4}}\]

\[\fitmath{A_{2}=\frac{\pi\cdot {(0.05)}^{2}}{4}\approx 0.001963495\,\mathrm{m^{2}}}\]

\item Aislar velocidad después de eliminar la velocidad de entrada.

\[\fitmath{V_{2}=\sqrt{\frac{2\cdot \mathrm{dp}}{\rho\cdot \left(1-{\left(\frac{A_{2}}{A_{1}}\right)}^{2}\right)}}}\]

\[\fitmath{V_{2}=\sqrt{\frac{2\cdot (15000)}{(998)\cdot \left(1-{\left(\frac{(0.001963495408)}{(0.007853981634)}\right)}^{2}\right)}}\approx 5.66252\,\mathrm{m/s}}\]

\item Convertir velocidad de garganta en caudal.

\[\fitmath{Q=A_{2}\cdot V_{2}}\]

\[\fitmath{Q=(0.001963495408)\cdot (5.662519603)\approx 0.01111833\,\mathrm{m^{3}/s}}\]

\end{enumerate}
